\begin{thebibliography}{99}
\footnotesize\itemsep=0pt

\bibitem{AndBrzRyb:ext}
Andruszkiewicz R.R., Brzezi\'nski T., Rybo{\l}owicz B., Ideal ring extensions
 and trusses, \href{https://doi.org/10.1016/j.jalgebra.2022.01.038}{\textit{J.~Algebra}} \textbf{600} (2022), 237--278,
 \href{https://arxiv.org/abs/2101.09484}{arXiv:2101.09484}.

\bibitem{5}
Baer R., Zur {E}inf\"uhrung des {S}charbegriffs, \href{https://doi.org/10.1515/crll.1929.160.199}{\textit{J.~Reine Angew. Math.}}
 \textbf{160} (1929), 199--207.

\bibitem{Ben:fib}
Benenti S., Fibr\'es affines canoniques et m\'ecanique newtonienne, in Action
 hamiltoniennes de groupes. {T}roisi\`eme th\'eor\`eme de {L}ie ({L}yon,
 1986), \textit{Travaux en Cours}, Vol.~27, Hermann, Paris, 1988, 13--37.

\bibitem{BreBrz:hea}
Breaz S., Brzezi\'nski T., Rybo{\l}owicz B., Saracco P., Heaps of modules and
 affine spaces, \textit{Ann. Mat. Pura Appl.}, {t}o appear, \href{https://arxiv.org/abs/2203.07268}{arXiv:2203.07268}.

\bibitem{3}
Brzezi\'nski T., Trusses: between braces and rings, \href{https://doi.org/10.1090/tran/7705}{\textit{Trans. Amer. Math.
 Soc.}} \textbf{372} (2019), 4149--4176, \href{https://arxiv.org/abs/1710.02870}{arXiv:1710.02870}.

\bibitem{Brz:par}
Brzezi\'nski T., Trusses: paragons, ideals and modules, \href{https://doi.org/10.1016/j.jpaa.2019.106258}{\textit{J.~Pure Appl.
 Algebra}} \textbf{224} (2020), 106258, 39~pages, \href{https://arxiv.org/abs/1901.07033}{arXiv:1901.07033}.

\bibitem{Brz:Lie}
Brzezi\'nski T., Lie trusses and heaps of Lie affebras, \href{https://doi.org/10.22323/1.406.0307}{\textit{Proc. Sci.}}
 \textbf{406} (2022), PoS(CORFU2021)307, 12~pages, \href{https://arxiv.org/abs/2203.12975}{arXiv:2203.12975}.

\bibitem{4}
Cari\~nena J.F., Grabowski J., Marmo G., Quantum bi-{H}amiltonian systems,
 \href{https://doi.org/10.1142/S0217751X00001954}{\textit{Internat.~J.~Modern Phys.~A}} \textbf{15} (2000), 4797--4810,
 \href{https://arxiv.org/abs/math-ph/0610011}{arXiv:math-ph/0610011}.

\bibitem{2}
Certaine J., The ternary operation {$(abc)=ab^{-1}c$} of a group, \href{https://doi.org/10.1090/S0002-9904-1943-08042-1}{\textit{Bull.
 Amer. Math. Soc.}} \textbf{49} (1943), 869--877.

\bibitem{7}
Gerstenhaber M., On the deformation of rings and algebras, \href{https://doi.org/10.2307/1970484}{\textit{Ann. of
 Math.}} \textbf{79} (1964), 59--103.

\bibitem{GraGra:Lie}
Grabowska K., Grabowski J., Urba\'nski P., Lie brackets on affine bundles,
 \href{https://doi.org/10.1023/A:1024457728027}{\textit{Ann. Global Anal. Geom.}} \textbf{24} (2003), 101--130,
 \href{https://arxiv.org/abs/math.DG/0203112}{arXiv:math.DG/0203112}.

\bibitem{Hoch}
Hochschild G., On the cohomology groups of an associative algebra, \href{https://doi.org/10.2307/1969145}{\textit{Ann.
 of Math.}} \textbf{46} (1945), 58--67.

\bibitem{Kos:Nij}
Kosmann-Schwarzbach Y., Nijenhuis structures on {C}ourant algebroids,
 \href{https://doi.org/10.1007/s00574-011-0032-5}{\textit{Bull. Braz. Math. Soc.~(N.S.)}} \textbf{42} (2011), 625--649,
 \href{https://arxiv.org/abs/1102.1410}{arXiv:1102.1410}.

\bibitem{Mag:sim}
Magri F., A simple model of the integrable {H}amiltonian equation,
 \href{https://doi.org/10.1063/1.523777}{\textit{J.~Math. Phys.}} \textbf{19} (1978), 1156--1162.

\bibitem{MasVig:non}
Massa E., Vignolo S., Bruno D., Non-holonomic {L}agrangian and {H}amiltonian
 mechanics: an intrinsic approach, \href{https://doi.org/10.1088/0305-4470/35/31/313}{\textit{J.~Phys.~A}} \textbf{35} (2002),
 6713--6742.

\bibitem{6}
Pr\"ufer H., Theorie der {A}belschen {G}ruppen, \href{https://doi.org/10.1007/BF01188079}{\textit{Math.~Z.}} \textbf{20}
 (1924), 165--187.

\bibitem{Tul:fra}
Tulczyjew W.M., Frame independence of analytical mechanics, \textit{Atti Accad.
 Sci. Torino Cl. Sci. Fis. Mat. Natur.} \textbf{119} (1985), 273--279.

\bibitem{Urb:aff}
Urba\'nski P., Affine {P}oisson structures in analytical mechanics, in
 Quantization and Infinite-Dimensional Systems ({B}ialowieza, 1993), \href{https://doi.org/10.1007/978-1-4615-2564-6_15}{Plenum},
 New York, 1994, 123--129.

\bibitem{Wei:uni}
Weinstein A., A universal phase space for particles in {Y}ang--{M}ills fields,
 \href{https://doi.org/10.1007/BF00400169}{\textit{Lett. Math. Phys.}} \textbf{2} (1978), 417--420.

\end{thebibliography}
